\documentclass[11pt,a4paper]{article}

\usepackage{fontspec}
\usepackage{polyglossia}

\setmainlanguage{russian}
\setotherlanguage{english}

\IfFontExistsTF{Times New Roman}
{\setmainfont{Times New Roman}}
{\setmainfont{Libertinus Serif}}

\IfFontExistsTF{Menlo}
{\newfontfamily\cyrillicfonttt{Menlo}}
{\newfontfamily\cyrillicfonttt{DejaVu Sans Mono}}

\usepackage[a4paper,
top=2cm,
bottom=2cm,
left=2.1cm,
right=2.1cm]{geometry}

\usepackage{amsmath,amssymb}
\usepackage{enumitem}
\usepackage{xurl}
\usepackage{hyperref}

\hypersetup{
    colorlinks=true,
    linkcolor=black,
    citecolor=black,
    urlcolor=blue
}

\setlength{\parindent}{1.25em}
\setlength{\parskip}{0.45em}
\setlength{\emergencystretch}{3em}
\linespread{1.04}\selectfont

\setlist[itemize]{
    topsep=3pt,
    itemsep=2pt,
    parsep=0pt,
    leftmargin=1.6em
}

\begin{document}

\begin{titlepage}
\centering

{\large Федеральное государственное автономное образовательное учреждение высшего образования\par}
{\large НАЦИОНАЛЬНЫЙ ИССЛЕДОВАТЕЛЬСКИЙ УНИВЕРСИТЕТ\par}
{\large «ВЫСШАЯ ШКОЛА ЭКОНОМИКИ»\par}
{\large Факультет экономических наук\par}
{\large Образовательная программа «Экономика»\par}

\vfill

{\Large\bfseries
Исследование устойчивости моделей волатильности и ценообразования опционов в различных рыночных режимах
\par}

\vfill

\begin{flushright}
\textbf{Выполнила:}\\
Ковалёва Мария Сергеевна\\
3 курс, группа БЭК224
\end{flushright}

\vspace{1cm}

\begin{flushright}
\textbf{Научный руководитель:}\\
Сизых Наталья Васильевна\\[0.5cm]
\underline{\hspace{5cm}}
\end{flushright}

\vfill

{\large Москва, 2026 г.\par}
\end{titlepage}

\tableofcontents
\newpage

\section{Введение}

\textbf{Актуальность исследования}


\textbf{Объект исследования}

Объектом исследования являются процессы ценообразования опционов на российском финансовом рынке

\textbf{Предмет исследования}

Предметом исследования выступает устойчивость моделей волатильности и ценообразования опционов в различных рыночных режимах, с целью повышения эффективность процесса оценки справедливой стоимости опционных контракто в зависимости от совокупность рыночных факторов.

\textbf{Проблема исследования.}

Классические модели определения цен опционов основаны на предположениях и допущениях. Они систематически отклоняются от наблюдаемых рыночных цен, которые мы принимаем за эталонные в данном исследовнии. 

\textbf{Гипотеза исследования.}

Отклонения модельных цен опционов от рыночных зависит от рыночного режима и ухудшается в периоды повышенной волатильности, что требует анализа и систематизации этих отклонений для дальнейшего учета этих зависимостей в построении моделей для оценки стоимости опционов. 

\textbf{Цель работы.}

Целью работы является исследование устойчивости моделей волатильности и ценообразования опционов в различных режимах волатильности, анализ ошибок моделей на реальных рыночных данных, разработка рекомендаций по корректному использованию исследуемых моделей в зависимости от рыночного состояния. 

\textbf{Задачи исследования.}

В соответствии с поставленной целью в работе решались следующие задачи:

\begin{itemize}

    \item 1. Изучить модели волатильности и ценообразования опционов, их предпосылки.
    \item 2. Сформировать и подготовить датасет рыночных цен опционов.
    \item 3. Реализовать базовые модели ценообразования опционов, расчитать модельные цены с их использованием. 
    \item 4. Провести сравнение модельных и рыночных цен опционов.
    \item 5. Проанализировать ошибки моделей в различных рыночных режимах.
    \item 6. Сравнить устойчивость моделей волатильности и подходов к моделированию цен опционов.

\end{itemize}

\textbf{Структура работы.}

Работа включает краткий теоретический обзор, методологию, эмпирическое исследование и список литературы.

\section{Анализ современных методов и моделей оценки волатильности и ценообразования опционов}

\subsection{Классические модели}

В работе рассматриваются классические модели ценообразования опционов, прежде всего биномиальная модель, модели Блэка--Шоулза и Блэка--76, как базовые инструменты для оценки стоимости опционов.

\subsection{Волатильность в ценообразовании опционов}

В работе кратко рассматриваются историческая, реализованная и подразумеваемая волатильность, а также улыбка волатильности, перекос волатильности и временная структура волатильности.

\subsection{Модели GARCH}

Отдельно рассматриваются модели GARCH как способ прогнозирования изменяющейся во времени волатильности базового актива.

\section{Методология исследования}

\subsection{Постановка исследовательской задачи}

В исследовании ставится задача построить единый рабочий датасет опционов и базового актива, рассчитать модельные цены по нескольким спецификациям модели Блэка--76, сравнить их с рыночными ценами и затем проверить, как меняются ошибки в зависимости от рыночного режима, срока до экспирации и денежности опциона.

\subsection{Сравнение модельных и рыночных цен}

В качестве рыночной цены опциона используется показатель SETTLEPRICE. Для каждого наблюдения дополнительно используются цена базового актива, страйк, срок до экспирации и тип опциона. Модельная цена рассчитывается в рамках модели Блэка--76 и затем сравнивается с наблюдаемой рыночной ценой.

Для каждого опциона рассчитывается ошибка:

\begin{equation}
error_i = P_i^{model} - P_i^{market}.
\end{equation}

Также рассчитываются абсолютная и относительная ошибки:

\begin{equation}
abs\_error_i = |P_i^{model} - P_i^{market}|,
\end{equation}

\begin{equation}
pct\_error_i =
\frac{
P_i^{model} - P_i^{market}
}{
P_i^{market}
}.
\end{equation}

\subsection{Метрики оценки качества моделей}

Для сравнения моделей используются метрики $MAE$, $RMSE$, средняя ошибка, медианная абсолютная ошибка и средняя абсолютная процентная ошибка. Основные формулы имеют вид:

\begin{equation}
MAE =
\frac{1}{N}
\sum_{i=1}^{N}
|P_i^{model} - P_i^{market}|,
\end{equation}

\begin{equation}
RMSE =
\sqrt{
\frac{1}{N}
\sum_{i=1}^{N}
(P_i^{model} - P_i^{market})^2
}.
\end{equation}

Затем агрегированные метрики рассчитываются как по полной выборке, так и по отдельным подвыборкам.

\subsection{Рыночные режимы}

Для исследования устойчивости моделей по рыночным состояниям формируются режимы низкой, средней и высокой волатильности. Дополнительно строятся разрезы по денежности опциона, сроку до экспирации и направлению движения базового актива.

\subsection{Исследовательский процесс}

Практическая часть выполняется в шесть последовательных этапов. На первом этапе проверяется качество исходного датасета и выбирается рабочее семейство контрактов. На втором этапе формируется итоговая выборка для моделирования: к опционам присоединяются дневные данные по базовому активу, рассчитываются срок до экспирации, денежность опциона и показатели волатильности. На третьем этапе строятся базовые модели Блэка--76 с использованием $hv\_21d$ и $hv\_63d$. На четвёртом этапе оценивается модель Блэка--76 с волатильностью по модели GARCH(1,1). На пятом этапе все модели сравниваются между собой по единым метрикам качества. На шестом этапе проводится формальная проверка устойчивости результатов с помощью регрессионного анализа и проверок устойчивости на подвыборках.

\section{Эмпирическое исследование устойчивости моделей}

\subsection{Описание данных}

На первом этапе через MOEX ISS API была собрана дневная история опционов за период с 3 января 2024 года по 19 мая 2026 года. После отбора рабочего семейства в выборке осталось 19574 наблюдения и 92 уникальных контрактов. Затем к опционным данным были присоединены дневные цены базового актива ИМОЕКС за период с 3 января 2024 года по 21 мая 2026 года. После мягкой фильтрации была сформирована итоговая выборка для моделирования объёмом 19327 наблюдений.

Для каждого наблюдения в рабочей выборке используются рыночная цена опциона, цена базового актива, страйк, срок до экспирации, тип опциона, историческая волатильность на окнах 21 и 63 дня, показатели доходности базового актива и индикаторы рыночного режима.

\subsection{Базовая схема ценообразования}

В качестве основной модели используется Блэка--76. Входные параметры задаются следующим образом: $F$ --- цена базового актива, $K$ --- страйк, $T$ --- срок до экспирации в годах, $r = 0.16$ --- постоянная безрисковая ставка, $\sigma$ --- выбранная оценка волатильности. Для оценки волатильности используются:
\begin{itemize}
    \item историческая волатильность на окне 21 день ($hv\_21d$);
    \item историческая волатильность на окне 63 дня ($hv\_63d$);
    \item прогноз волатильности на основе модели GARCH(1,1).
\end{itemize}

\subsection{Расчёт модельных цен}

Для каждой из трёх спецификаций рассчитываются модельные цены по формуле:

\begin{equation}
P^{model} = Black76(F, K, T, r, \sigma),
\end{equation}

после чего для каждого опциона вычисляются $error$, $abs\_error$, $pct\_error$ и $squared\_error$.

\subsection{Анализ ошибок моделей}

Анализ ошибок выполняется по рыночным режимам, денежности опциона и сроку до экспирации. Для модели Блэка--76 с $hv\_21d$ на полной выборке получены значения $MAE \approx 326.29$ и $RMSE \approx 452.59$. Для модели Блэка--76 с $hv\_63d$ получены значения $MAE \approx 309.00$ и $RMSE \approx 435.45$. Для модели Блэка--76 с волатильностью по модели GARCH получены значения $MAE \approx 326.79$, $RMSE \approx 455.04$, $mean\_error \approx -309.25$ и $median\_abs\_error \approx 215.58$.

По рыночным режимам для лучшей базовой модели Блэка--76 с $hv\_63d$ получены следующие значения $MAE$: в режиме низкой волатильности --- около $170.86$, в режиме средней волатильности --- около $357.62$, в режиме высокой волатильности --- около $411.14$. По срокам до экспирации для той же модели получены значения $MAE$: для интервала $0$--$7$ дней --- около $10.74$, для интервала $8$--$30$ дней --- около $16.95$, для интервала $31$--$90$ дней --- около $50.40$, для интервала свыше $91$ дня --- около $356.48$.

\subsection{Прогнозирование волатильности с помощью GARCH}

Для логарифмических доходностей базового актива оценивается модель GARCH(1,1):

\begin{equation}
\sigma_t^2 =
\omega
+
\alpha \varepsilon_{t-1}^2
+
\beta \sigma_{t-1}^2.
\end{equation}

После оценки модели условная дневная волатильность переводится в годовую по формуле $\sigma^{annual}_t = \sqrt{252}\sigma_t$, затем присоединяется к опционной выборке по торговой дате и используется как входной параметр в модели Блэка--76. Для этой спецификации получены агрегированные метрики: $N = 19291$, $MAE \approx 326.79$, $RMSE \approx 455.04$, $mean\_error \approx -309.25$, $median\_abs\_error \approx 215.58$.

\subsection{Сравнение моделей}

На этапе итогового сравнения сопоставляются три спецификации:
\begin{itemize}
    \item Блэка--76 с $hv\_21d$;
    \item Блэка--76 с $hv\_63d$;
    \item Блэка--76 с волатильностью по модели GARCH.
\end{itemize}

Полученные результаты:
\begin{itemize}
    \item для модели Блэка--76 с $hv\_21d$: $MAE \approx 326.29$, $RMSE \approx 452.59$;
    \item для модели Блэка--76 с $hv\_63d$: $MAE \approx 309.00$, $RMSE \approx 435.45$;
    \item для модели Блэка--76 с волатильностью по модели GARCH: $MAE \approx 326.79$, $RMSE \approx 455.04$, $mean\_error \approx -309.25$, $median\_abs\_error \approx 215.58$.
\end{itemize}

Минимальные значения ошибок на полной выборке получены у модели Блэка--76 с $hv\_63d$. В сводной таблице сравнения для неё зафиксированы наименьшие значения как $MAE$, так и $RMSE$.

\subsection{Эконометрическая проверка результатов и проверки устойчивости}

На заключительном этапе для абсолютной ошибки лучшей базовой модели, то есть для $abs\_error\_{hv\_63d}$, оцениваются три линейные регрессии с робастными стандартными ошибками типа $HC3$. Сначала оценивается однофакторная модель только с индикатором высокого режима волатильности, затем добавляются $DTE$ и абсолютная логарифмическая денежность, а в расширенной спецификации дополнительно учитываются направление тренда и тип опциона.

Формально оцениваются регрессии вида:

\begin{equation}
|e_i| = \alpha + \beta_1 HighVol_i + u_i,
\end{equation}

\begin{equation}
|e_i| = \alpha + \beta_1 HighVol_i + \beta_2 DTE_i + \beta_3 |\log(M_i)| + u_i.
\end{equation}

В однофакторной спецификации коэффициент при переменной высокого режима волатильности составил примерно $150.69$ и оказался статистически значимым. Во второй спецификации коэффициент при $DTE$ составил около $0.7635$, коэффициент при абсолютной логарифмической денежности --- около $118.87$, а коэффициент при индикаторе высокого режима волатильности перестал быть статистически значимым. В расширенной спецификации коэффициент при высоком режиме волатильности составил около $13.72$ и оказался значимым, коэффициент при $DTE$ --- около $0.6338$, коэффициент при абсолютной логарифмической денежности --- около $121.96$.

Проверки устойчивости выполняются для полной выборки, для выборки без крайних хвостов денежности и для краткосрочных опционов с $DTE \leq 91$. Для модели Блэка--76 с $hv\_63d$ на полной выборке получены $MAE \approx 309.00$ и $RMSE \approx 435.45$, для выборки без крайних хвостов денежности --- $MAE \approx 269.08$ и $RMSE \approx 378.49$, для краткосрочных опционов --- $MAE \approx 43.09$ и $RMSE \approx 59.69$. Для модели GARCH на краткосрочной подвыборке получены $MAE \approx 41.50$ и $RMSE \approx 58.73$.

\begin{thebibliography}{99}

\bibitem{Hull}
Hull J. C. Options, Futures, and Other Derivatives. --- 10th ed. --- Pearson, 2018.

\bibitem{Black}
Black F. The Pricing of Commodity Contracts // Journal of Financial Economics. --- 1976. --- Vol. 3. --- P. 167--179.

\bibitem{BlackScholes}
Black F., Scholes M. The Pricing of Options and Corporate Liabilities // Journal of Political Economy. --- 1973. --- Vol. 81, No. 3. --- P. 637--654.

\bibitem{Merton}
Merton R. Theory of Rational Option Pricing // Bell Journal of Economics and Management Science. --- 1973. --- Vol. 4, No. 1. --- P. 141--183.

\bibitem{Bollerslev}
Bollerslev T. Generalized Autoregressive Conditional Heteroskedasticity // Journal of Econometrics. --- 1986. --- Vol. 31. --- P. 307--327.

\bibitem{Engle}
Engle R. Autoregressive Conditional Heteroscedasticity with Estimates of the Variance of United Kingdom Inflation // Econometrica. --- 1982. --- Vol. 50, No. 4. --- P. 987--1007.

\bibitem{Heston}
Heston S. A Closed-Form Solution for Options with Stochastic Volatility with Applications to Bond and Currency Options // Review of Financial Studies. --- 1993. --- Vol. 6, No. 2. --- P. 327--343.

\bibitem{Bakshi}
Bakshi G., Cao C., Chen Z. Empirical Performance of Alternative Option Pricing Models // Journal of Finance. --- 1997. --- Vol. 52, No. 5. --- P. 2003--2049.

\bibitem{Dumas}
Dumas B., Fleming J., Whaley R. Implied Volatility Functions: Empirical Tests // Journal of Finance. --- 1998. --- Vol. 53, No. 6. --- P. 2059--2106.

\bibitem{Christoffersen}
Christoffersen P., Jacobs K. The Importance of the Loss Function in Option Valuation // Journal of Financial Economics. --- 2004. --- Vol. 72. --- P. 291--318.

\bibitem{Cont}
Cont R. Empirical Properties of Asset Returns: Stylized Facts and Statistical Issues // Quantitative Finance. --- 2001. --- Vol. 1, No. 2. --- P. 223--236.

\bibitem{Tsay}
Tsay R. Analysis of Financial Time Series. --- 3rd ed. --- Wiley, 2010.

\bibitem{Shreve}
Shreve S. Stochastic Calculus for Finance II: Continuous-Time Models. --- Springer, 2004.

\bibitem{Haug}
Haug E. G. The Complete Guide to Option Pricing Formulas. --- McGraw-Hill, 2007.

\bibitem{Gatheral}
Gatheral J. The Volatility Surface: A Practitioner's Guide. --- Wiley, 2006.

\end{thebibliography}

\end{document}
